% ============================================================
\chapter{Simulation and Queueing Theory}
\label{ch:simulation}
% ============================================================

\section{Introduction}

Simulation is an essential tool in operations research when analytical models
are too complex or intractable. This chapter covers discrete event simulation,
random variate generation, Monte Carlo estimation, and variance reduction
techniques.

\begin{intuition}
When a system is too complex for mathematical analysis (queues with priorities,
logistics networks, etc.), we \emph{simulate} it: generate random realizations
of the system and observe the results statistically. It is like running
thousands of virtual experiments.
\end{intuition}

\section{Discrete event simulation}

\begin{definition}[Discrete event simulation (DES)]
\index{discrete event simulation}
\textbf{Discrete event simulation} models a system as a sequence of
instantaneous events (customer arrival, service completion, failure, etc.)
that modify the system state at discrete time points.
\end{definition}

\begin{definition}[DES components]
\begin{itemize}
  \item \textbf{System state}: variables describing the system at time $t$
        (number of customers, server status, etc.).
  \item \textbf{Simulation clock}: current time $t$.
  \item \textbf{Future event list} (FEL): list of scheduled events, sorted
        by time.
  \item \textbf{Statistical counters}: accumulators for performance measures.
\end{itemize}
\end{definition}

\begin{algorithme}[title=\textbf{DES main loop}]
\begin{enumerate}
  \item Initialize state, clock $t = 0$, and FEL.
  \item While the stopping condition is not met:
    \begin{enumerate}
      \item Extract the next event $(t_e, \text{type})$ from the FEL.
      \item Advance the clock: $t \leftarrow t_e$.
      \item Process the event (update state, schedule new events).
      \item Update statistical counters.
    \end{enumerate}
  \item Compute and return performance indicators.
\end{enumerate}
\end{algorithme}

\section{Random variate generation}

\begin{theorem}[Inverse transform method]
\index{inverse transform method}
If $U \sim \mathcal{U}(0,1)$ and $F$ is the CDF of a continuous random
variable $X$, then:
\[
  X = F^{-1}(U)
\]
has the distribution of $X$.
\end{theorem}

\begin{proof}
For any $x \in \R$:
$\PP(F^{-1}(U) \leq x) = \PP(U \leq F(x)) = F(x)$,
since $U \sim \mathcal{U}(0,1)$ and $F$ is nondecreasing.
\end{proof}

\begin{example}[Exponential distribution]
For $X \sim \mathrm{Exp}(\lambda)$, $F(x) = 1 - e^{-\lambda x}$, so
$F^{-1}(u) = -\frac{1}{\lambda} \ln(1-u)$. We generate:
\[
  X = -\frac{1}{\lambda} \ln(U), \quad U \sim \mathcal{U}(0,1).
\]
(We use $\ln(U)$ instead of $\ln(1-U)$ since $U$ and $1-U$ have the same
distribution.)
\end{example}

\begin{keyformulas}[title=\textbf{Generating classical distributions}]
\begin{itemize}
  \item \textbf{Exponential}$(\lambda)$:
        $X = -\frac{1}{\lambda} \ln(U)$.
  \item \textbf{Bernoulli}$(p)$: $X = \ind_{U \leq p}$.
  \item \textbf{Poisson}$(\lambda)$: count $U_i$ until
        $\prod U_i < e^{-\lambda}$.
  \item \textbf{Normal} (Box-Muller):
        $X = \sqrt{-2 \ln U_1} \cos(2\pi U_2)$.
\end{itemize}
\end{keyformulas}

\begin{theorem}[Acceptance-rejection method]
\index{acceptance-rejection}
Let $f$ be the target density and $g$ a proposal density with
$f(x) \leq M g(x)$ for all $x$. To generate $X \sim f$:
\begin{enumerate}
  \item Generate $Y \sim g$ and $U \sim \mathcal{U}(0,1)$.
  \item If $U \leq \frac{f(Y)}{M g(Y)}$, accept $X = Y$; otherwise reject
        and repeat.
\end{enumerate}
The acceptance rate is $1/M$.
\end{theorem}

\section{Monte Carlo estimation}

\begin{definition}[Monte Carlo estimator]
\index{Monte Carlo estimation}
To estimate $\theta = \E[h(X)]$, the Monte Carlo estimator is:
\[
  \hat{\theta}_n = \frac{1}{n} \sum_{i=1}^n h(X_i),
\]
where $X_1, \ldots, X_n$ are i.i.d.\ copies of $X$.
\end{definition}

\begin{theorem}[Properties of the MC estimator]
\begin{enumerate}
  \item \textbf{Unbiased}: $\E[\hat{\theta}_n] = \theta$.
  \item \textbf{Consistency}: $\hat{\theta}_n \xrightarrow{\text{a.s.}} \theta$
        (strong law of large numbers).
  \item \textbf{CLT}: $\sqrt{n}(\hat{\theta}_n - \theta)
        \xrightarrow{d} \mathcal{N}(0, \sigma^2)$ where
        $\sigma^2 = \Var(h(X))$.
  \item \textbf{95\% confidence interval}:
        $\hat{\theta}_n \pm 1.96 \cdot \hat{\sigma}/\sqrt{n}$.
\end{enumerate}
\end{theorem}

\begin{example}[Estimating $\pi$]
Generate points $(X, Y)$ uniformly in $[0,1]^2$ and compute the proportion
$\hat{p}$ satisfying $X^2 + Y^2 \leq 1$. Then $\hat{\pi} = 4\hat{p}$.
\end{example}

\section{Variance reduction techniques}

\begin{definition}[Antithetic variates]
\index{antithetic variates}
Instead of $n$ independent samples $U_1, \ldots, U_n$, use the pairs
$(U_i, 1 - U_i)$. If $h$ is monotone, the negative covariance between
$h(U)$ and $h(1-U)$ reduces the variance of the mean.
\end{definition}

\begin{proposition}[Antithetic variance reduction]
If $h$ is monotone and $U \sim \mathcal{U}(0,1)$:
\[
  \Var\!\left(\frac{h(U) + h(1-U)}{2}\right)
  \leq \frac{\Var(h(U))}{2}.
\]
\end{proposition}

\begin{definition}[Control variates]
\index{control variates}
Let $Y$ be an auxiliary variable with known mean $\E[Y] = \mu_Y$. The
control variate estimator is:
\[
  \hat{\theta}_c = \hat{\theta}_n - c(\bar{Y}_n - \mu_Y),
\]
with optimal coefficient $c^* = \Cov(h(X), Y) / \Var(Y)$.
\end{definition}

\begin{definition}[Importance sampling]
\index{importance sampling}
To estimate $\theta = \E_f[h(X)]$, sample from a proposal distribution $g$
and reweight:
\[
  \hat{\theta}_{\mathrm{IS}} = \frac{1}{n} \sum_{i=1}^n h(X_i)
  \frac{f(X_i)}{g(X_i)}, \quad X_i \sim g.
\]
The optimal proposal is $g^*(x) \propto \abs{h(x)} f(x)$.
\end{definition}

\section{Applications in operations research}

\begin{itemize}
  \item \textbf{Queueing systems}: simulation of M/M/1, M/G/k with
        priorities and breakdowns.
  \item \textbf{Inventory management}: $(s, S)$ policy with random demand.
  \item \textbf{Logistics}: supply chain simulation.
  \item \textbf{Finance}: option pricing via Monte Carlo.
  \item \textbf{Reliability}: estimating failure probability of complex
        systems.
\end{itemize}

\section{Python implementation}

\begin{codeblock}[title=\textbf{M/M/1 discrete event simulation}]
\begin{minted}{python}
import numpy as np
import heapq

def simulate_mm1(lam, mu, n_clients):
    """Simulate an M/M/1 queue."""
    t = 0.0
    fel = []  # future event list (time, type)
    n_in_system = 0
    arrivals = 0
    total_wait = 0.0
    heapq.heappush(fel,
        (np.random.exponential(1/lam), 'A'))
    while arrivals < n_clients:
        event_time, event_type = heapq.heappop(fel)
        t = event_time
        if event_type == 'A':  # Arrival
            arrivals += 1
            n_in_system += 1
            if n_in_system == 1:
                service = np.random.exponential(1/mu)
                heapq.heappush(fel, (t + service, 'D'))
            heapq.heappush(fel,
                (t + np.random.exponential(1/lam), 'A'))
        else:  # Departure
            n_in_system -= 1
            if n_in_system > 0:
                service = np.random.exponential(1/mu)
                heapq.heappush(fel, (t + service, 'D'))
    return t / arrivals
\end{minted}
\end{codeblock}

\section{Exercises}

\begin{exercise}[$\star$ -- Inverse transform]
Generate 1000 realizations of an exponential distribution with parameter
$\lambda = 2$ using the inverse transform method. Compare the histogram to
the theoretical density.
\end{exercise}

\begin{exercise}[$\star\star$ -- M/M/1 simulation]
Simulate an M/M/1 queue with $\lambda = 4$ and $\mu = 5$ for 10,000
customers. Estimate the mean waiting time and compare to the theoretical
value $W = \frac{1}{\mu - \lambda}$.
\end{exercise}

\begin{exercise}[$\star\star$ -- Monte Carlo integration]
Estimate $I = \int_0^1 e^{-x^2}\,dx$ by Monte Carlo with $n = 10^4$.
Construct a 95\% confidence interval. Compare with antithetic variates.
\end{exercise}

\begin{exercise}[$\star\star\star$ -- Control variates]
Estimate $\E[e^U]$ where $U \sim \mathcal{U}(0,1)$. Use $Y = U$ as a control
variate. Compute the optimal coefficient $c^*$ and the theoretical variance
reduction.
\end{exercise}

\begin{exercise}[$\star\star\star$ -- Acceptance-rejection]
Generate realizations of a $\mathrm{Beta}(2,5)$ distribution by
acceptance-rejection with $g = \mathcal{U}(0,1)$. Compute the constant $M$
and the acceptance rate.
\end{exercise}
